\documentclass[11pt,a4paper]{article}

\usepackage[utf8]{inputenc}
\usepackage[T1]{fontenc}
\usepackage[spanish]{babel}
\usepackage{amsmath,amssymb,mathtools}
\usepackage{geometry}
\usepackage{lmodern}
\usepackage{microtype}

\geometry{margin=2.4cm}
\setlength{\parindent}{0pt}
\setlength{\parskip}{0.45em}

\newcommand{\R}{\mathbb{R}}
\newcommand{\sen}{\operatorname{sen}}
\newcommand{\rg}{\operatorname{rg}}
\newcommand{\tr}{\operatorname{tr}}

\begin{document}

\begin{center}
{\LARGE\bfseries Geometría Lineal}\\[0.4em]
{\Large Tema 4: Movimientos afines}\\[0.4em]
{\large Soluciones}
\end{center}

\section{Ejercicio 1}
\subsection{Solución}

La recta es:

\[
r=(1,2,3)+\mathcal L\{(1,1,1)\}.
\]

Por tanto:

\[
\boxed{
\begin{cases}
x=1+\lambda,\\
y=2+\lambda,\\
z=3+\lambda.
\end{cases}
}
\]

Eliminando el parámetro:

\[
\boxed{
x-1=y-2=z-3
}
\]

\section{Ejercicio 2}
\subsection{Solución}

Tomamos:

\[
P_1=(1,1,1,0),
\quad
P_2=(0,1,2,3),
\quad
P_3=(0,0,1,-1).
\]

Los vectores directores son:

\[
P_1P_2=(-1,0,1,3),
\]

\[
P_1P_3=(-1,-1,0,-1).
\]

Dos vectores normales al plano son:

\[
n_1=(1,-1,1,0),
\qquad
n_2=(3,-4,0,1).
\]

Por tanto:

\[
\boxed{
\begin{cases}
x-y+z-1=0,\\
3x-4y+t+1=0.
\end{cases}
}
\]

\section{Ejercicio 3}
\subsection{Solución}

Por la relación de Chasles:

\[
P_1P_2+P_2P_3=P_1P_3.
\]

Continuando:

\[
P_1P_2+P_2P_3+\cdots+P_{n-1}P_n
=
P_1P_n.
\]

Además:

\[
P_nP_1=-P_1P_n.
\]

Por tanto:

\[
\boxed{
P_1P_2+P_2P_3+\cdots+P_{n-1}P_n+P_nP_1=0
}
\]

\section{Ejercicio 4}
\subsection{Solución}

Sean:

\[
r=P+\mathcal L\{u\},
\qquad
s=Q+\mathcal L\{v\}.
\]

Si:

\[
\mathcal L\{u\}=\mathcal L\{v\},
\]

las rectas son paralelas.

Si no son paralelas, \(u\) y \(v\) son linealmente independientes y:

\[
\mathcal L\{u,v\}=\R^2.
\]

Por tanto:

\[
PQ\in\mathcal L\{u,v\},
\]

y existe una única solución de:

\[
P+\lambda u=Q+\mu v.
\]

Así:

\[
\boxed{
\text{En }\mathbb A^2\text{ dos rectas o son paralelas o se cortan en un punto.}
}
\]

\section{Ejercicio 5}
\subsection{Solución}

Sea:

\[
r=P+\mathcal L\{u\}
\]

y:

\[
\pi=Q+S,
\qquad
\dim S=2.
\]

Si:

\[
u\in S,
\]

la recta es paralela al plano.

Si:

\[
u\notin S,
\]

entonces:

\[
S+\mathcal L\{u\}=\R^3.
\]

Por tanto, existe una única solución de:

\[
P+\lambda u\in\pi.
\]

Así:

\[
\boxed{
\text{En }\mathbb A^3\text{ una recta y un plano o son paralelos o se cortan en un punto.}
}
\]

\section{Ejercicio 6}
\subsection{Solución}

Se tiene:

\[
P=(3,4),
\qquad
O^*=(1,3).
\]

Por tanto:

\[
O^*P=(2,1).
\]

Buscamos:

\[
(2,1)
=
x^*(1,2)
+
y^*(-2,1).
\]

Así:

\[
\begin{cases}
x^*-2y^*=2,\\
2x^*+y^*=1.
\end{cases}
\]

Resolviendo:

\[
x^*=\frac45,
\qquad
y^*=-\frac35.
\]

Por tanto:

\[
\boxed{
P=
\left(
\frac45,-\frac35
\right)_{R^*}
}
\]

\section{Ejercicio 7}
\subsection{Solución}

Un conjunto de cuatro puntos de \(\mathbb A^3\) es un sistema de referencia afín si:

\[
P_0P_1,\quad P_0P_2,\quad P_0P_3
\]

son linealmente independientes.

Para el apartado a):

\[
\det
\begin{pmatrix}
0&1&1\\
-1&-1&0\\
0&-1&0
\end{pmatrix}
=1\neq0.
\]

Por tanto:

\[
\boxed{
\text{a) Sí es un sistema de referencia afín.}
}
\]

Para el apartado b), todos los vectores tienen tercera coordenada nula, por lo que:

\[
\det=0.
\]

Así:

\[
\boxed{
\text{b) No es un sistema de referencia afín.}
}
\]

Para el apartado c):

\[
\det
\begin{pmatrix}
1&1&-1\\
0&1&-1\\
2&0&0
\end{pmatrix}
=0.
\]

Por tanto:

\[
\boxed{
\text{c) No es un sistema de referencia afín.}
}
\]

\section{Ejercicio 8}
\subsection{Solución}

Las nuevas bases verifican:

\[
v'=v-2w,
\qquad
w'=v+w.
\]

Por tanto:

\[
\begin{pmatrix}
x\\
y
\end{pmatrix}_{R}
=
\begin{pmatrix}
1&1\\
-2&1
\end{pmatrix}
\begin{pmatrix}
x'\\
y'
\end{pmatrix}_{R^*}.
\]

La matriz inversa es:

\[
\frac13
\begin{pmatrix}
1&-1\\
2&1
\end{pmatrix}.
\]

Así:

\[
\boxed{
\begin{pmatrix}
x'\\
y'
\end{pmatrix}_{R^*}
=
\frac13
\begin{pmatrix}
1&-1\\
2&1
\end{pmatrix}
\begin{pmatrix}
x\\
y
\end{pmatrix}_{R}
}
\]

\section{Ejercicio 9}
\subsection{Solución}

El nuevo sistema es:

\[
R^*
=
\{O^*;-e_1+e_3,-e_2+e_3,e_1+e_2+e_3\},
\]

con:

\[
O^*=(1,1,0).
\]

Si:

\[
X=(x^*,y^*,z^*)_{R^*},
\]

entonces, en el sistema \(R\):

\[
\begin{cases}
x=1-x^*+z^*,\\
y=1-y^*+z^*,\\
z=x^*+y^*+z^*.
\end{cases}
\]

Sustituyendo en:

\[
x+y+z=2,
\]

se obtiene:

\[
2+3z^*=2.
\]

Por tanto:

\[
\boxed{
z^*=0
}
\]

\section{Ejercicio 10}
\subsection{Solución}

Sea \(f\) una aplicación afín y \(\bar f\) su aplicación vectorial asociada.

Si \(f\) es inyectiva y:

\[
\bar f(v)=0,
\]

tomamos puntos \(P,Q\) tales que:

\[
PQ=v.
\]

Entonces:

\[
f(P)f(Q)=0,
\]

por lo que:

\[
f(P)=f(Q).
\]

Como \(f\) es inyectiva:

\[
P=Q
\]

y:

\[
v=0.
\]

Así, \(\bar f\) es inyectiva.

Recíprocamente, si \(\bar f\) es inyectiva y:

\[
f(P)=f(Q),
\]

entonces:

\[
\bar f(PQ)=f(P)f(Q)=0.
\]

Por tanto:

\[
PQ=0
\]

y:

\[
P=Q.
\]

Así:

\[
\boxed{
f\text{ es inyectiva}\iff\bar f\text{ es inyectiva}
}
\]

\section{Ejercicio 11}
\subsection{Solución}

Si \(f\) es suprayectiva, dados \(u\in V'\), tomamos puntos \(P',Q'\) tales que:

\[
P'Q'=u.
\]

Existen \(P,Q\) con:

\[
f(P)=P',
\qquad
f(Q)=Q'.
\]

Por tanto:

\[
\bar f(PQ)
=
f(P)f(Q)
=
u.
\]

Así, \(\bar f\) es suprayectiva.

Recíprocamente, si \(\bar f\) es suprayectiva, fijamos \(P\in A\).

Para cualquier \(Q'\in A'\), existe \(v\in V\) tal que:

\[
\bar f(v)=f(P)Q'.
\]

Existe un único punto:

\[
Q=P+v.
\]

Entonces:

\[
f(Q)
=
f(P)+\bar f(v)
=
Q'.
\]

Por tanto:

\[
\boxed{
f\text{ es suprayectiva}\iff\bar f\text{ es suprayectiva}
}
\]

\section{Ejercicio 12}
\subsection{Solución}

La parte lineal de la transformación es:

\[
A=
\begin{pmatrix}
1&3&4\\
2&1&3\\
-1&2&1
\end{pmatrix}.
\]

Su determinante es:

\[
\det A=0.
\]

Por tanto, la aplicación lineal asociada no es biyectiva.

Así:

\[
\boxed{
\text{La transformación afín no es biyectiva.}
}
\]

\section{Ejercicio 13}
\subsection{Solución}

Sea una traslación:

\[
f(P)=P+v,
\qquad
v\neq0.
\]

Si \(P\) fuese fijo:

\[
P+v=P.
\]

Entonces:

\[
v=0,
\]

lo que contradice la hipótesis.

Por tanto:

\[
\boxed{
\text{Una traslación distinta de la identidad no tiene puntos fijos.}
}
\]

\section{Ejercicio 14}
\subsection{Solución}

Sean \(P\) y \(Q\) dos puntos fijos:

\[
f(P)=P,
\qquad
f(Q)=Q.
\]

Todo punto de la recta \(PQ\) puede escribirse como:

\[
X=P+\lambda PQ.
\]

Como \(f\) es afín:

\[
f(X)
=
f(P)+\lambda\bar f(PQ).
\]

Además:

\[
\bar f(PQ)
=
f(P)f(Q)
=
PQ.
\]

Por tanto:

\[
f(X)
=
P+\lambda PQ
=
X.
\]

Así:

\[
\boxed{
\text{Toda la recta }PQ\text{ queda fija por }f.
}
\]

\section{Ejercicio 15}
\subsection{Solución}

Sean \(P,Q,R\) tres puntos fijos no alineados.

Los vectores:

\[
PQ,\qquad PR
\]

son linealmente independientes y forman una base de \(\R^2\).

Como:

\[
\bar f(PQ)=PQ,
\qquad
\bar f(PR)=PR,
\]

la aplicación vectorial asociada es la identidad.

Además:

\[
f(P)=P.
\]

Por tanto:

\[
\boxed{
f=\operatorname{Id}_{\mathbb A^2}
}
\]

\section{Ejercicio 16}
\subsection{Solución}

La transformación es:

\[
(x',y')=(4x-3y,x).
\]

Los puntos fijos satisfacen:

\[
\begin{cases}
4x-3y=x,\\
x=y.
\end{cases}
\]

Ambas ecuaciones equivalen a:

\[
x=y.
\]

Por tanto:

\[
\boxed{
\{P\text{ fijo}\}
=
\{(x,y):x=y\}
}
\]

Además:

\[
PP'
=
(x'-x,y'-y).
\]

Así:

\[
PP'
=
(3x-3y,x-y)
=
(x-y)(3,1).
\]

Por tanto:

\[
\boxed{
PP'\parallel(3,1)
}
\]

\section{Ejercicio 17}
\subsection{Solución}

La transformación es:

\[
X'
=
\begin{pmatrix}
3\\
1
\end{pmatrix}
+
\begin{pmatrix}
-\frac35&-\frac45\\
-\frac45&\frac35
\end{pmatrix}
X.
\]

La parte lineal es una simetría respecto de una recta.

No existen puntos fijos, por lo que se trata de una simetría deslizante.

El origen se transforma en:

\[
O'=(3,1).
\]

El punto medio es:

\[
P=
\left(
\frac32,\frac12
\right).
\]

Su imagen es:

\[
P'=
\left(
\frac{17}{10},
\frac1{10}
\right).
\]

Así:

\[
PP'
=
\left(
\frac15,-\frac25
\right).
\]

La recta invariante pasa por \(P\) y tiene dirección \((1,-2)\):

\[
\boxed{
4x+2y-7=0
}
\]

El vector de desplazamiento es:

\[
\boxed{
v=
\left(
\frac15,-\frac25
\right)
}
\]

\section{Ejercicio 18}
\subsection{Solución}

Una homotecia de centro \(C\) y razón \(2\) cumple:

\[
A'=C+2(A-C).
\]

Por tanto:

\[
C=2A-A'.
\]

Sustituyendo:

\[
C
=
2(1,1,2)-(2,-2,6).
\]

Así:

\[
\boxed{
C=(0,4,-2)
}
\]

\section{Ejercicio 19}
\subsection{Solución}

La primera homotecia es:

\[
H_{O,3}(X)=3X.
\]

Como:

\[
A=(0,2),
\]

su imagen es:

\[
A'=(0,6).
\]

La segunda homotecia es:

\[
H_{A',2}(X)
=
A'+2(X-A')
=
2X-A'.
\]

Por tanto:

\[
H_{A',2}\circ H_{O,3}(X)
=
6X-A'.
\]

Así:

\[
f(X)
=
6X+(0,-6).
\]

Una homotecia de razón \(6\) cumple:

\[
f(X)=6X+(1-6)C.
\]

Por tanto:

\[
-5C=(0,-6).
\]

Así:

\[
\boxed{
\text{La composición es la homotecia de razón }6
\text{ y centro }
C=
\left(
0,\frac65
\right).
}
\]

\section{Ejercicio 20}
\subsection{Solución}

La transformación tiene la forma:

\[
X'=B+AX.
\]

Como:

\[
O=(0,0)
\longmapsto
O'=(15,-1),
\]

se tiene:

\[
B=
\begin{pmatrix}
15\\
-1
\end{pmatrix}.
\]

Además:

\[
A
\begin{pmatrix}
1\\
2
\end{pmatrix}
=
\begin{pmatrix}
-22\\
-4
\end{pmatrix},
\]

\[
A
\begin{pmatrix}
2\\
1
\end{pmatrix}
=
\begin{pmatrix}
-20\\
10
\end{pmatrix}.
\]

Resolviendo:

\[
A=
\begin{pmatrix}
-6&-8\\
8&-6
\end{pmatrix}.
\]

Por tanto:

\[
\boxed{
\begin{pmatrix}
x'\\
y'
\end{pmatrix}
=
\begin{pmatrix}
15\\
-1
\end{pmatrix}
+
\begin{pmatrix}
-6&-8\\
8&-6
\end{pmatrix}
\begin{pmatrix}
x\\
y
\end{pmatrix}
}
\]

Se verifica:

\[
AA^t=100I.
\]

Por tanto:

\[
\boxed{
\text{Sí es una semejanza, de razón }10.
}
\]

\section{Ejercicio 21}
\subsection{Solución}

La homotecia de centro:

\[
C=(0,0,1)
\]

y razón \(2\) es:

\[
H(x,y,z)
=
(2x,2y,2z-1).
\]

Aplicando después:

\[
f(x,y,z)
=
(2x+y-1,3z-2,4x+y),
\]

se obtiene:

\[
(f\circ H)(x,y,z)
=
(4x+2y-1,6z-5,8x+2y).
\]

Por tanto:

\[
\boxed{
\begin{pmatrix}
x'\\
y'\\
z'
\end{pmatrix}
=
\begin{pmatrix}
-1\\
-5\\
0
\end{pmatrix}
+
\begin{pmatrix}
4&2&0\\
0&0&6\\
8&2&0
\end{pmatrix}
\begin{pmatrix}
x\\
y\\
z
\end{pmatrix}
}
\]

La matriz lineal no verifica:

\[
AA^t=k^2I.
\]

Por tanto:

\[
\boxed{
\text{La transformación resultante no es una semejanza.}
}
\]

\section{Ejercicio 22}
\subsection{Solución}

La transformación es:

\[
\begin{pmatrix}
x'\\
y'
\end{pmatrix}
=
\begin{pmatrix}
1\\
-1
\end{pmatrix}
+
I
\begin{pmatrix}
x+1\\
y+4
\end{pmatrix}.
\]

Así:

\[
(x',y')=(x+2,y+3).
\]

Por tanto:

\[
\boxed{
\text{Es una traslación de vector }(2,3).
}
\]

\section{Ejercicio 23}
\subsection{Solución}

La transformación es:

\[
X'
=
\begin{pmatrix}
2\\
2
\end{pmatrix}
+
\begin{pmatrix}
0&-1\\
-1&0
\end{pmatrix}
X.
\]

La matriz lineal es una simetría.

Los puntos fijos satisfacen:

\[
\begin{cases}
x=2-y,\\
y=2-x.
\end{cases}
\]

Por tanto:

\[
x+y=2.
\]

Así:

\[
\boxed{
\text{Es la simetría respecto de la recta }x+y=2.
}
\]

\section{Ejercicio 24}
\subsection{Solución}

La recta es:

\[
x-y-1=0.
\]

Un vector normal es:

\[
n=(1,-1).
\]

La reflexión viene dada por:

\[
X'
=
X
-
2\frac{x-y-1}{2}(1,-1).
\]

Así:

\[
x'=y+1,
\qquad
y'=x-1.
\]

Por tanto:

\[
\boxed{
\begin{pmatrix}
x'\\
y'
\end{pmatrix}
=
\begin{pmatrix}
1\\
-1
\end{pmatrix}
+
\begin{pmatrix}
0&1\\
1&0
\end{pmatrix}
\begin{pmatrix}
x\\
y
\end{pmatrix}
}
\]

\section{Ejercicio 25}
\subsection{Solución}

La transformación es:

\[
\begin{pmatrix}
x'\\
y'
\end{pmatrix}
=
\begin{pmatrix}
1\\
0
\end{pmatrix}
+
\begin{pmatrix}
0&1\\
1&0
\end{pmatrix}
\begin{pmatrix}
x\\
y+1
\end{pmatrix}.
\]

Por tanto:

\[
(x',y')=(y+2,x).
\]

No tiene puntos fijos.

La parte lineal es una simetría, por lo que se trata de una simetría deslizante.

El origen se transforma en:

\[
O'=(2,0).
\]

Su punto medio es:

\[
P=(1,0).
\]

La imagen de \(P\) es:

\[
P'=(2,1).
\]

Así:

\[
PP'=(1,1).
\]

Por tanto, el eje es:

\[
\boxed{
y=x-1
}
\]

y el desplazamiento:

\[
\boxed{
v=(1,1).
}
\]

\section{Ejercicio 26}
\subsection{Solución}

La recta de simetría es:

\[
y=2x-1,
\]

equivalentemente:

\[
2x-y-1=0.
\]

La matriz de la reflexión es:

\[
R=
\frac15
\begin{pmatrix}
-3&4\\
4&3
\end{pmatrix}.
\]

La parte independiente de la reflexión es:

\[
b_R=
\begin{pmatrix}
\frac45\\[1mm]
-\frac25
\end{pmatrix}.
\]

Añadiendo el desplazamiento:

\[
v=(-1,-2),
\]

se obtiene:

\[
b=
\begin{pmatrix}
-\frac15\\[1mm]
-\frac{12}{5}
\end{pmatrix}.
\]

Por tanto:

\[
\boxed{
\begin{pmatrix}
x'\\
y'
\end{pmatrix}
=
\begin{pmatrix}
-\frac15\\[1mm]
-\frac{12}{5}
\end{pmatrix}
+
\frac15
\begin{pmatrix}
-3&4\\
4&3
\end{pmatrix}
\begin{pmatrix}
x\\
y
\end{pmatrix}
}
\]

\section{Ejercicio 27}
\subsection{Solución}

La transformación es:

\[
X'
=
\begin{pmatrix}
\frac15\\[1mm]
\frac25
\end{pmatrix}
+
\begin{pmatrix}
-\frac35&-\frac45\\[1mm]
\frac45&-\frac35
\end{pmatrix}
\left(
X-
\begin{pmatrix}
\frac15\\[1mm]
\frac25
\end{pmatrix}
\right).
\]

Tiene un único punto fijo:

\[
\boxed{
C=
\left(
\frac15,\frac25
\right)
}
\]

La matriz tiene determinante \(1\), luego es un giro.

Además:

\[
\cos\theta=-\frac35,
\qquad
\sen\theta=\frac45.
\]

Por tanto:

\[
\boxed{
\theta
=
\arccos\left(-\frac35\right)
}
\]

en sentido positivo.

\section{Ejercicio 28}
\subsection{Solución}

El giro tiene centro:

\[
C=(1,-1)
\]

y ángulo:

\[
\theta=45^\circ.
\]

Su matriz lineal es:

\[
R=
\frac1{\sqrt2}
\begin{pmatrix}
1&-1\\
1&1
\end{pmatrix}.
\]

Por tanto:

\[
X'=C+R(X-C).
\]

Así:

\[
\boxed{
\begin{pmatrix}
x'\\
y'
\end{pmatrix}
=
\begin{pmatrix}
1\\
-1
\end{pmatrix}
+
\frac1{\sqrt2}
\begin{pmatrix}
1&-1\\
1&1
\end{pmatrix}
\left[
\begin{pmatrix}
x\\
y
\end{pmatrix}
-
\begin{pmatrix}
1\\
-1
\end{pmatrix}
\right]
}
\]

\section{Ejercicio 29}
\subsection{Solución}

Consideramos:

\[
G_1\circ G_2.
\]

Los ángulos son:

\[
\theta_1=\frac{\pi}{3},
\qquad
\theta_2=\frac{\pi}{6}.
\]

Por tanto, el ángulo de la composición es:

\[
\theta
=
\frac{\pi}{2}.
\]

La matriz lineal es:

\[
R=
\begin{pmatrix}
0&-1\\
1&0
\end{pmatrix}.
\]

La parte independiente resulta:

\[
b=
\begin{pmatrix}
-\frac12+\frac{3\sqrt3}{2}\\[1mm]
\frac12-\frac{\sqrt3}{2}
\end{pmatrix}.
\]

El centro satisface:

\[
(I-R)C=b.
\]

Resolviendo:

\[
\boxed{
C=
\left(
-\frac12+\sqrt3,
\frac{\sqrt3}{2}
\right)
}
\]

Por tanto:

\[
\boxed{
G_1\circ G_2
\text{ es un giro de ángulo }
\frac{\pi}{2}
\text{ y centro }
\left(
-\frac12+\sqrt3,
\frac{\sqrt3}{2}
\right).
}
\]

\section{Ejercicio 30}
\subsection{Solución}

La parte lineal es:

\[
A=
\begin{pmatrix}
1&b\\
a-1&a^2
\end{pmatrix}.
\]

Para que sea un movimiento rígido:

\[
A^tA=I.
\]

La primera columna debe tener norma \(1\):

\[
1+(a-1)^2=1.
\]

Por tanto:

\[
a=1.
\]

La ortogonalidad entre las columnas da:

\[
b=0.
\]

Así:

\[
\boxed{
a=1,
\qquad
b=0,
\qquad
c\in\R.
}
\]

La transformación queda:

\[
(x',y')=(x+c,y).
\]

Por tanto:

\[
\boxed{
\begin{cases}
c=0 & \text{identidad},\\
c\neq0 & \text{traslación de vector }(c,0).
\end{cases}
}
\]

\section{Ejercicio 31}
\subsection{Solución}

La ecuación matricial equivale a:

\[
\begin{pmatrix}
x'\\
y'\\
z'
\end{pmatrix}
=
\begin{pmatrix}
1\\
2\\
3
\end{pmatrix}
+
I
\begin{pmatrix}
x\\
y\\
z
\end{pmatrix}.
\]

Por tanto:

\[
\boxed{
\text{Es una traslación de vector }(1,2,3).
}
\]

\section{Ejercicio 32}
\subsection{Solución}

La transformación es:

\[
X'=B+AX
\]

con:

\[
B=
\begin{pmatrix}
-\frac{10}{3}\\[1mm]
-\frac{20}{3}\\[1mm]
-\frac{10}{3}
\end{pmatrix}
\]

y:

\[
A=
\frac13
\begin{pmatrix}
2&-2&-1\\
-2&-1&-2\\
-1&-2&2
\end{pmatrix}.
\]

Se verifica:

\[
A^tA=I,
\qquad
\det A=-1.
\]

Los autovalores son:

\[
1,\quad1,\quad-1.
\]

El subespacio asociado a \(-1\) es:

\[
V_{-1}
=
\mathcal L\{(1,2,1)\}.
\]

Los puntos fijos satisfacen:

\[
(A-I)X=-B.
\]

Se obtiene:

\[
x+2y+z+10=0.
\]

Por tanto:

\[
\boxed{
\text{Es la simetría especular respecto del plano }
x+2y+z+10=0.
}
\]

\section{Ejercicio 33}
\subsection{Solución}

El plano es:

\[
x-2y=7.
\]

Tomamos:

\[
n=(1,-2,0).
\]

La simetría es:

\[
X'
=
X
-
2\frac{x-2y-7}{5}(1,-2,0).
\]

Por tanto:

\[
\boxed{
\begin{pmatrix}
x'\\
y'\\
z'
\end{pmatrix}
=
\begin{pmatrix}
\frac{14}{5}\\[1mm]
-\frac{28}{5}\\[1mm]
0
\end{pmatrix}
+
\frac15
\begin{pmatrix}
3&4&0\\
4&-3&0\\
0&0&5
\end{pmatrix}
\begin{pmatrix}
x\\
y\\
z
\end{pmatrix}
}
\]

\section{Ejercicio 34}
\subsection{Solución}

La transformación es:

\[
X'=B+AX,
\]

con:

\[
B=
\begin{pmatrix}
-6\\
1\\
-7
\end{pmatrix},
\qquad
A=
\begin{pmatrix}
0&0&-1\\
0&1&0\\
-1&0&0
\end{pmatrix}.
\]

Se verifica:

\[
A^tA=I,
\qquad
\det A=-1.
\]

Además:

\[
V_{-1}
=
\mathcal L\{(1,0,1)\}.
\]

No existen puntos fijos, por lo que es una simetría deslizante respecto de un plano.

Descomponemos:

\[
B=B_\parallel+B_\perp,
\]

obteniendo:

\[
B_\parallel
=
\left(
\frac12,1,-\frac12
\right),
\]

\[
B_\perp
=
\left(
-\frac{13}{2},0,-\frac{13}{2}
\right).
\]

El plano de simetría es:

\[
\boxed{
x+z=-\frac{13}{2}
}
\]

y el vector de desplazamiento:

\[
\boxed{
v=
\left(
\frac12,1,-\frac12
\right).
}
\]

\section{Ejercicio 35}
\subsection{Solución}

El plano es:

\[
x+y+1=0.
\]

La simetría respecto de dicho plano es:

\[
\begin{pmatrix}
x_r\\
y_r\\
z_r
\end{pmatrix}
=
\begin{pmatrix}
-1\\
-1\\
0
\end{pmatrix}
+
\begin{pmatrix}
0&-1&0\\
-1&0&0\\
0&0&1
\end{pmatrix}
\begin{pmatrix}
x\\
y\\
z
\end{pmatrix}.
\]

Añadiendo el desplazamiento:

\[
v=(0,0,2),
\]

se obtiene:

\[
\boxed{
\begin{pmatrix}
x'\\
y'\\
z'
\end{pmatrix}
=
\begin{pmatrix}
-1\\
-1\\
2
\end{pmatrix}
+
\begin{pmatrix}
0&-1&0\\
-1&0&0\\
0&0&1
\end{pmatrix}
\begin{pmatrix}
x\\
y\\
z
\end{pmatrix}
}
\]

\section{Ejercicio 36}
\subsection{Solución}

La transformación tiene la forma:

\[
X'
=
C+A(X-C),
\]

con:

\[
C=(0,0,1)
\]

y:

\[
A=
\frac13
\begin{pmatrix}
1&1+\sqrt3&1-\sqrt3\\
1-\sqrt3&1&1+\sqrt3\\
1+\sqrt3&1-\sqrt3&1
\end{pmatrix}.
\]

Se verifica:

\[
A^tA=I,
\qquad
\det A=1.
\]

Además:

\[
A(1,1,1)=(1,1,1).
\]

Por tanto, el eje es:

\[
\boxed{
r:
X=(0,0,1)+\lambda(1,1,1)
}
\]

La traza es:

\[
\tr A=1.
\]

Para un giro en \(\R^3\):

\[
\tr A=1+2\cos\theta.
\]

Así:

\[
\cos\theta=0.
\]

Por tanto:

\[
\boxed{
\theta=\frac{\pi}{2}
}
\]

y la transformación es un giro de \(90^\circ\) alrededor de \(r\).

\section{Ejercicio 37}
\subsection{Solución}

El eje es:

\[
r:
(1,0,0)+\lambda(-1,-1,1).
\]

Tomamos como dirección orientada:

\[
u=
\frac1{\sqrt3}(-1,-1,1).
\]

Para:

\[
\theta=240^\circ,
\]

la fórmula de Rodrigues proporciona:

\[
R=
\begin{pmatrix}
0&1&0\\
0&0&-1\\
-1&0&0
\end{pmatrix}.
\]

Con:

\[
C=(1,0,0),
\]

el giro es:

\[
X'=C+R(X-C).
\]

Por tanto:

\[
\boxed{
\begin{pmatrix}
x'\\
y'\\
z'
\end{pmatrix}
=
\begin{pmatrix}
1\\
0\\
1
\end{pmatrix}
+
\begin{pmatrix}
0&1&0\\
0&0&-1\\
-1&0&0
\end{pmatrix}
\begin{pmatrix}
x\\
y\\
z
\end{pmatrix}
}
\]

\section{Ejercicio 38}
\subsection{Solución}

La transformación es:

\[
X'=B+AX,
\]

con:

\[
B=(1,1,0)
\]

y:

\[
A=
\begin{pmatrix}
0&1&0\\
0&0&-1\\
-1&0&0
\end{pmatrix}.
\]

Se verifica:

\[
A^tA=I,
\qquad
\det A=1.
\]

Además:

\[
V_1
=
\mathcal L\{(-1,-1,1)\}.
\]

No existen puntos fijos, por lo que es un movimiento helicoidal.

Como:

\[
\tr A=0,
\]

se tiene:

\[
1+2\cos\theta=0.
\]

Por tanto:

\[
\boxed{
\theta=120^\circ
}
\]

La componente de \(B\) paralela al eje es:

\[
\boxed{
v=
\left(
\frac23,\frac23,-\frac23
\right).
}
\]

Un punto del eje es:

\[
C=
\left(
\frac23,\frac13,0
\right).
\]

Así:

\[
\boxed{
r:
X=
\left(
\frac23,\frac13,0
\right)
+
\lambda(-1,-1,1)
}
\]

\section{Ejercicio 39}
\subsection{Solución}

El eje pasa por:

\[
C=(1,1,1)
\]

y tiene dirección:

\[
v=(0,1,1).
\]

El giro es de:

\[
\theta=\frac{\pi}{2}
\]

en sentido \(-v\).

Tomamos:

\[
u=
\frac1{\sqrt2}(0,-1,-1).
\]

Por Rodrigues:

\[
R=
\begin{pmatrix}
0&\frac{\sqrt2}{2}&-\frac{\sqrt2}{2}\\[1mm]
-\frac{\sqrt2}{2}&\frac12&\frac12\\[1mm]
\frac{\sqrt2}{2}&\frac12&\frac12
\end{pmatrix}.
\]

El movimiento helicoidal es:

\[
X'=C+R(X-C)+v.
\]

Por tanto:

\[
\boxed{
X'
=
\begin{pmatrix}
1\\[1mm]
1+\frac{\sqrt2}{2}\\[1mm]
1-\frac{\sqrt2}{2}
\end{pmatrix}
+
\begin{pmatrix}
0&\frac{\sqrt2}{2}&-\frac{\sqrt2}{2}\\[1mm]
-\frac{\sqrt2}{2}&\frac12&\frac12\\[1mm]
\frac{\sqrt2}{2}&\frac12&\frac12
\end{pmatrix}
X
}
\]

\section{Ejercicio 40}
\subsection{Solución}

La transformación es:

\[
X'=B+AX,
\]

con:

\[
B=
\begin{pmatrix}
\frac13\\[1mm]
\frac43\\[1mm]
-\frac23
\end{pmatrix}
\]

y:

\[
A=
\frac13
\begin{pmatrix}
-2&-2&-1\\
-2&1&2\\
1&-2&2
\end{pmatrix}.
\]

Se verifica:

\[
A^tA=I,
\qquad
\det A=-1.
\]

Tiene un único punto fijo:

\[
\boxed{
C=
\left(
\frac12,0,-\frac32
\right).
}
\]

Además:

\[
V_{-1}
=
\mathcal L\{(2,1,0)\}.
\]

Por tanto, es una simetría rotatoria.

Su traza es:

\[
\tr A=\frac13.
\]

Así:

\[
-1+2\cos\theta=\frac13,
\]

por lo que:

\[
\boxed{
\theta
=
\arccos\left(\frac23\right)
}
\]

El eje es:

\[
\boxed{
r:
X=
\left(
\frac12,0,-\frac32
\right)
+
\lambda(2,1,0)
}
\]

y el plano de simetría:

\[
\boxed{
2x+y-1=0.
}
\]

\section{Ejercicio 41}
\subsection{Solución}

La simetría respecto del plano:

\[
x=1
\]

es:

\[
\boxed{
\begin{pmatrix}
x'\\
y'\\
z'
\end{pmatrix}
=
\begin{pmatrix}
2\\
0\\
0
\end{pmatrix}
+
\begin{pmatrix}
-1&0&0\\
0&1&0\\
0&0&1
\end{pmatrix}
\begin{pmatrix}
x\\
y\\
z
\end{pmatrix}
}
\]

La recta perpendicular al plano que pasa por:

\[
(0,0,1)
\]

es:

\[
r:
(0,0,1)+\lambda(1,0,0).
\]

El giro de \(90^\circ\) alrededor de dicha recta es:

\[
\boxed{
\begin{pmatrix}
x'\\
y'\\
z'
\end{pmatrix}
=
\begin{pmatrix}
0\\
1\\
1
\end{pmatrix}
+
\begin{pmatrix}
1&0&0\\
0&0&-1\\
0&1&0
\end{pmatrix}
\begin{pmatrix}
x\\
y\\
z
\end{pmatrix}
}
\]

\section{Ejercicio 42}
\subsection{Solución}

La parte lineal es:

\[
A=
\begin{pmatrix}
1&0&0\\
0&\alpha&\frac{\sqrt3}{2}\\
0&\frac{\sqrt3}{2}&\frac12
\end{pmatrix}.
\]

Para que sea ortogonal:

\[
A^tA=I.
\]

De la ortogonalidad de las columnas segunda y tercera:

\[
\frac{\sqrt3}{2}\alpha
+
\frac{\sqrt3}{4}
=
0.
\]

Por tanto:

\[
\boxed{
\alpha=-\frac12
}
\]

Para este valor:

\[
\det A=-1,
\qquad
\tr A=1.
\]

La transformación no tiene puntos fijos, por lo que es una simetría deslizante respecto de un plano.

El subespacio:

\[
V_{-1}
=
\mathcal L\{(0,-\sqrt3,1)\}.
\]

El plano de simetría es:

\[
\boxed{
-\sqrt3\,y+z=-\frac{\sqrt3}{2}
}
\]

y el vector de desplazamiento:

\[
\boxed{
v=
\left(
2,\frac14,\frac{\sqrt3}{4}
\right).
}
\]

\section{Ejercicio 43}
\subsection{Solución}

Los planos son:

\[
\pi_1:
3y-4z+6=0,
\]

\[
\pi_2:
3y-4z-4=0.
\]

Son paralelos y tienen vector normal:

\[
n=(0,3,-4).
\]

La composición de las simetrías respecto de dos planos paralelos es una traslación perpendicular a ambos.

Tomando la composición en el orden:

\[
S_{\pi_2}\circ S_{\pi_1},
\]

el vector de traslación es:

\[
v
=
\frac{2(4-(-6))}{25}(0,3,-4).
\]

Por tanto:

\[
\boxed{
v=
\left(
0,\frac{12}{5},-\frac{16}{5}
\right)
}
\]

Así:

\[
\boxed{
X'
=
X+
\left(
0,\frac{12}{5},-\frac{16}{5}
\right)
}
\]

y la transformación resultante es una traslación.

\section{Ejercicio 44}
\subsection{Solución}

El plano es:

\[
3x-4z+5=0.
\]

Su vector normal es:

\[
n=(3,0,-4).
\]

La matriz lineal de la simetría es:

\[
R=
\frac1{25}
\begin{pmatrix}
7&0&24\\
0&25&0\\
24&0&-7
\end{pmatrix}.
\]

La parte independiente es:

\[
b_S
=
\left(
-\frac65,0,\frac85
\right).
\]

La traslación tiene vector:

\[
v=(15,0,-20).
\]

Tomando:

\[
T_v\circ S,
\]

se obtiene:

\[
b=b_S+v
=
\left(
\frac{69}{5},0,-\frac{92}{5}
\right).
\]

Por tanto:

\[
\boxed{
X'
=
\begin{pmatrix}
\frac{69}{5}\\[1mm]
0\\[1mm]
-\frac{92}{5}
\end{pmatrix}
+
\frac1{25}
\begin{pmatrix}
7&0&24\\
0&25&0\\
24&0&-7
\end{pmatrix}
X
}
\]

Como el desplazamiento es perpendicular al plano, la composición sigue siendo una simetría especular.

Su plano de puntos fijos es:

\[
\boxed{
3x-4z-\frac{115}{2}=0
}
\]

Por tanto:

\[
\boxed{
\text{La transformación resultante es una simetría respecto de un plano.}
}
\]

\end{document}
